% 2-6-xai.tex  | Budget 2pp -- CURRENTLY THIN / PARTIALLY BLOCKED
% SOURCE MAP: interpretability named as open challenge ->oconnor_review_2025 ;
%  interpretability-by-design (coefficient bounds) ->ghelasi_day-ahead_2025 ;
%  dedicated XAI+DL for electricity prices ->trebbien_explainable_2023 (TITLE ONLY).
% !!! trebbien_explainable_2023 has NO abstract in references.bib. Only a
%     title-level mention is grounded. To reach a full 2pp and detail the actual
%     XAI methods (SHAP / feature attribution) this section should cover, PASTE a
%     passage from trebbien_2023. lucas_price_2020 is NOT clearly XAI from its
%     abstract, so it is NOT cited here (moved to 2-7, UK balancing market).
\section{تفسیرپذیری و هوش مصنوعی توضیح‌پذیر}

با پیچیده‌تر و دقیق‌ترشدنِ مدل‌های یادگیری ماشین و یادگیری عمیق، پرسشی تازه اهمیت
یافته است: چرا مدل یک پیش‌بینیِ خاص را ارائه می‌دهد؟ مرورهای اخیر، تفسیرپذیریِ
مدل‌ها را در کنارِ کیفیتِ داده و تعمیم‌پذیری، از جمله چالش‌های همچنان گشوده‌ی این
حوزه برمی‌شمارند \cite{oconnor_review_2025}. بسیاری از مدل‌های پرکاربرد \lr{—}
به‌ویژه معماری‌های عمیق \lr{—} ماهیتی جعبه‌سیاه دارند و بی‌آن‌که سازوکارِ درونیِ
آن‌ها روشن باشد، صرفاً بر پایه‌ی دقتِ خروجی ارزیابی می‌شوند.

در پاسخ به این نیاز، دو رویکردِ کلی در ادبیات دیده می‌شود. رویکردِ نخست، تعبیه‌ی
تفسیرپذیری در خودِ ساختارِ مدل است؛ برای نمونه، مدلی اقتصادسنجی با محدودکردنِ
ضرایب به بازه‌هایی که از مدل‌های بنیادی استخراج می‌شوند، هم دقت و هم تفسیرپذیری را
هم‌زمان دنبال می‌کند \cite{ghelasi_day-ahead_2025}. رویکردِ دوم، به‌کارگیریِ
روش‌های هوشِ مصنوعیِ توضیح‌پذیر به‌صورتِ پسینی برای تحلیلِ مدل‌های پیچیده است؛
پژوهشی به‌طورِ اختصاصی هوشِ مصنوعیِ توضیح‌پذیر و یادگیری عمیق را برای تحلیل و
پیش‌بینیِ سری‌های زمانیِ پیچیده، از جمله قیمتِ برق، به‌کار گرفته است
\cite{trebbien_explainable_2023}.

با این‌همه، در قیاس با انبوهِ مطالعاتِ متمرکز بر دقت، شمارِ کارهایی که به‌طورِ
نظام‌مند به تفسیرپذیریِ مدل‌های پیش‌بینیِ قیمت می‌پردازند بسیار اندک است. همین
کم‌توجهی، یکی از شکاف‌های اصلیِ این حوزه و از انگیزه‌های محوریِ پژوهشِ حاضر برای
به‌کارگیریِ تحلیلِ تفسیرپذیری در کنارِ سنجشِ دقت است.
